\section{Additional results}
\label{app:results}

\subsection{Downstream clinical utility (atypia classifier)}
\label{app:clinical}

\textbf{Protocol.} The proposal's clinical-utility test trains a classifier on raw Aperio $\times$20 frames
of the eleven MITOS-ATYPIA-14 training slides to predict the adjudicated nuclear-atypia score (1--3), then
evaluates it on held-out Hamamatsu frames before and after normalisation. The recovery is the accuracy on
normalised frames minus the accuracy on raw Hamamatsu frames, so a positive value means normalisation
helped the classifier.

\textbf{A retracted result.} The first classifier (a ResNet18, \texttt{atypia\_r18}) appeared to favour
diffusion: the A3 rung at strength 0.50 gave a recovery of $+0.0625$, against $+0.0457$ for the best
classical method, Reinhard. A later audit found that every method had been run \AtoH{}, producing
Hamamatsu-styled images, whereas an Aperio-trained classifier needs \HtoA{} outputs. An A0 output was twice
as close in pixel error to its Aperio source as to its Hamamatsu reference, which confirmed the direction
error. The claim was retracted. Re-running in the correct direction, no configuration beats the raw input:
our source-conditioned \HtoA{} model gives $-0.025$ (CI crossing zero), and with F3 fusion the recovery is
exactly $0.0$ (95\% CI $[-0.025,+0.025]$, 120 frames), even though fusion raises its SSIM from 0.548 to
0.730. The rerun also found and fixed a second bug, in which the raw-Hamamatsu reference silently read
Aperio pixels.

\textbf{The classifier itself is not valid.} \texttt{atypia\_r18} predicts score 2 for about 91\% of held-out
frames even on raw Aperio, its own domain, where it scores 0.55 against 0.50 for always answering 2. Its
0.947 validation accuracy came from two slides containing 46 score-2 frames out of 47, so it is the
majority rate rather than skill. The labels are close to slide-determined: 19 of the 23 score-1 training
frames come from slide A12, and all 38 held-out score-1 frames come from A13 and A16. A retrained variant
with augmentation and class-balanced sampling was worse (raw-Hamamatsu accuracy 0.442 to 0.192).

\textbf{Two rebuilt classifiers, both gated.} Both use frozen UNI2-h features, nested leave-one-slide-out
cross-validation on the 297 training frames, and five pre-registered gates: balanced accuracy $>0.45$,
quadratic weighted kappa (QWK) $>0.2$, no class predicted for more than 80\% of frames, a slide-bootstrap
lower bound on balanced accuracy above chance ($1/3$), and non-zero recall for score 1 on the held-out-in-CV
slide A12. Table~\ref{tab:atypia} gives the results. The linear probe is at chance. Of three v2 candidates
(an ordinal head, attention-based multiple-instance pooling, and a two-stage model that first predicts the
x40 atypia criteria), the criteria model clears four of five gates but fails the bootstrap lower bound:
its interval $[0.26, 0.73]$ cannot be separated from chance with only eleven slides, and the threshold was
not relaxed afterwards. A separate within-slide test does pass. With each slide's mean removed from both
predictions and labels, so that recognising the slide cannot help, the predicted criteria correlate with
the rated ones (mean Spearman 0.240, CI $[0.148, 0.318]$; nuclei size 0.378, anisonucleosis 0.299,
nucleoli size 0.408, nuclei contour $-0.125$). Frozen features therefore carry real atypia signal, but not
yet enough for a valid score classifier. By protocol the held-out slides were never scored with these
models.

\begin{table}[!htbp]
\centering
\caption{Rebuilt atypia classifiers, nested leave-one-slide-out cross-validation on 297 training frames
(always predicting score 2 gives accuracy 0.747). Brackets: slide-bootstrap 95\% intervals.}
\label{tab:atypia}
\small
\begin{tabular}{lrrrcc}
\toprule
Model & Accuracy & Balanced accuracy & QWK & Recall (1 / 2 / 3) & Gates passed\\
\midrule
UNI2-h logistic probe & 0.579 & 0.312 [0.23, 0.46] & 0.022 & 0.00 / 0.73 / 0.21 & 1 of 5\\
v2 M1: ordinal head & 0.633 & 0.331 [0.27, 0.47] & 0.105 & 0.00 / 0.80 / 0.19 & 0 of 5\\
v2 M2: attention MIL & 0.478 & 0.292 [0.19, 0.50] & 0.098 & 0.00 / 0.57 / 0.31 & 1 of 5\\
v2 M3: criteria two-stage & 0.525 & 0.614 [0.26, 0.73] & 0.304 & 0.83 / 0.50 / 0.52 & 4 of 5\\
\bottomrule
\end{tabular}
\end{table}

\subsection{Other closed negatives}
\textbf{Differential diffusion.} To protect structure inside the diffusion process, each pixel was given
its own denoising strength from a Canny-derived change map, low near edges and high elsewhere. Six
configurations (dilation radius $\in\{2,4,8\}$, maximum strength $\in\{0.50,0.70\}$) gave SSIM 0.162--0.167
on 20 internal crops, against $0.162\pm0.006$ for the plain model, so there was no gain
(Fig.~\ref{fig:app-diffdiff}). On dense H\&E the protected and free regions interleave finely, and the
shared U-Net receptive field plausibly lets generation bleed into protected pixels. This failure is what
motivated fusion outside the U-Net.

\textbf{Round-trip consistency.} Translating A06 crops \AtoH{} and back to Aperio (50-step DDIM, strength
0.20) gives SSIM 0.143, PSNR 13.1\,dB and LAB distance 28.2 against the original
over 64 crops: the round trip does not return the input (Fig.~\ref{fig:app-roundtrip}).

\subsection{SDXL diagnostics and the decision gate}
\textbf{Decision gate.} At the end of Phase~1 the ladder configuration A4 was chosen as the base for SDXL,
and the A5 warm-start transfer was also run because A5 beat A4 on A06 by 2.6--8.7 in LAB distance, far above
the 0.1--0.3 run-to-run noise floor.

\textbf{Why native-1024 SDXL fails at 50 pairs.} Five diagnostics were run on the 50-pair 1024\,px model.
(D1) A baseline mismatch was ruled out: re-measured against a do-nothing baseline computed natively at
1024\,px, the model still gives $\dlab=-5.25$. (D2) Removing the LoRA makes colour worse ($-8.21$ against $-5.25$), so the LoRA helps by about 3
LAB but is too weak. (D3) RGB-only conditioning matches RGB plus Canny, and Canny alone collapses (SSIM 0.16,
$\dlab=-17.9$), so the edge channels are inert. (D4) No ControlNet scale gives positive recovery. (D5)
Raising the LoRA scale from 1.0 to 2.0 makes recovery monotonically worse ($-4.38$ to $-8.00$). The model is
data-starved, which the 96- and 144-pair runs confirm.

\textbf{The \AtoH{}/\HtoA{} asymmetry.} SD\,1.5 fails \HtoA{} while SDXL succeeds. Two diagnostics compared
the mean LAB of model outputs with the real scanner means (domain gap 17.97). Zero-source ControlNet
outputs refuted the idea that the asymmetry follows the backbone's default colour, because img2img leaks the
source colour into them. Pure text-to-image samples (Fig.~\ref{fig:app-t2i}) gave partial support (five of
six comparisons consistent), and SDXL at 1024\,px \HtoA{} remains an open counterexample: its default colour
leans the wrong way, yet it recovers $\dlab=+9.20$.

\clearpage
